% =====================================================================
\section{Le volet décisionnel}

Le volet décisionnel part des \textbf{mêmes tables} que l'application. Rien
n'est saisi deux fois : trois chemins de restitution s'appuient sur une même
base opérationnelle.

\begin{center}
\small
\begin{tabularx}{\textwidth}{|>{\columncolor{btkpale}\bfseries}L{3.6cm}|X|}
\hline
Tableau de bord embarqué & Servlets et services de l'application, graphiques ApexCharts --- temps réel, dans la page\\
\hline
Power BI & Vues \texttt{V\_BI\_*} lues directement par Power BI Desktop --- rapport de quatre pages\\
\hline
Segmentation & Datamart produit par la chaîne ETL, puis apprentissage non supervisé en Python\\
\hline
\end{tabularx}
\end{center}

\subsection{Les vues décisionnelles}

Quatre vues forment la \textbf{couche sémantique} : elles traduisent les
colonnes techniques en notions métier, de sorte que Power BI n'ait ni jointure
ni calcul compliqué à refaire.

\noindent\begin{tabularx}{\textwidth}{|>{\ttfamily}L{3.5cm}|L{3.2cm}|X|}
\hline
\textmd{\textbf{Vue}} & \textbf{Fichier} & \textbf{Contenu}\\
\hline
V\_BI\_EMPLOYES & v\_bi\_employes.sql & effectifs, rôle calculé, statut\\
\hline
V\_BI\_CLIENTS & v\_bi\_clients.sql & portefeuille : identité, sexe, type, statut\\
\hline
V\_BI\_OBJECTIFS & v\_bi\_objectifs.sql & production : comptes, crédits, cartes, épargne\\
\hline
V\_BI\_POINTAGE & setup\_pointage.sql & présence : indicateurs présent / retard / absent\\
\hline
\end{tabularx}

\vspace{0.6em}
\texttt{V\_BI\_OBJECTIFS} illustre le principe : les colonnes brutes de
\texttt{B\_OBJECTIF} sont consolidées et renommées en indicateurs, et les clés
sont remplacées par des libellés.

\begin{lstlisting}[style=sql]
CREATE OR REPLACE VIEW V_BI_OBJECTIFS AS
SELECT
    O.ID_OBJECTIF, O.ANNEE_MOIS,
    A.LIBELLE_AGENCE      AS NOM_AGENCE,
    U.LIBELLE_UTILISATEUR AS NOM_GESTIONNAIRE,
    O.SITUATION_ADMINISTRATIVE,
    (NVL(O.SOUSCRIPTION_COMPTE_CHEQUES_OA, 0)
   + NVL(O.SOUSCRIPTION_COMPTE_EPARGNES_OA, 0)
   + NVL(O.SOUSCRIPTION_COMPTE_COURANTS_OA, 0)) AS TOTAL_OUVERTURE_COMPTES,
    NVL(O.PRODUCTION_CREDITS_CONSO_OA, 0)          AS CREDITS_CONSO,
    NVL(O.PRODUCTION_CREDITS_IMMO_OA, 0)           AS CREDITS_IMMO,
    NVL(O.PRODUCTION_CREDITS_INVESTISSEMENT_OA, 0) AS CREDITS_INVEST,
    NVL(O.SOUSCRIPTION_CARTES_UNITE_OA, 0)         AS CARTES_COMMANDES,
    NVL(O.EPARGNE_ADD_OA, 0)                       AS COLLECTE_EPARGNE
FROM B_OBJECTIF O
LEFT JOIN AGENCE         A ON O.SK_AGENCE      = A.SK_AGENCE
LEFT JOIN B_UTILISATEURS U ON O.SK_UTILISATEUR = U.SK_UTILISATEUR;
\end{lstlisting}

\begin{retenir}
Le \texttt{NVL(\dots, 0)} systématique et les jointures \texttt{LEFT JOIN} sont
volontaires : une agence sans gestionnaire renseigné, ou un indicateur non
saisi, ne doit pas faire disparaître la ligne du rapport. C'est la première
règle d'une couche sémantique : \textbf{ne jamais perdre de lignes en chemin}.
\end{retenir}

\subsection{Le schéma \texttt{BTK\_BI}}

Power BI ne se connecte pas au schéma applicatif. Le script
\texttt{powerbi/creer\_schema\_bi.sql} crée un utilisateur et un schéma dédiés,
puis y recopie les tables nécessaires :

\begin{lstlisting}[style=sql]
CREATE USER BTK_BI IDENTIFIED BY "...";
GRANT CONNECT, RESOURCE TO BTK_BI;
ALTER USER BTK_BI QUOTA UNLIMITED ON USERS;

CREATE TABLE BTK_BI.AGENCE         AS SELECT * FROM SYSTEM.AGENCE;
CREATE TABLE BTK_BI.B_UTILISATEURS AS SELECT * FROM SYSTEM.B_UTILISATEURS;
CREATE TABLE BTK_BI.B_CLIENTS      AS SELECT * FROM SYSTEM.B_CLIENTS;
CREATE TABLE BTK_BI.B_OBJECTIF     AS SELECT * FROM SYSTEM.B_OBJECTIF;
CREATE TABLE BTK_BI.POINTAGE       AS SELECT * FROM SYSTEM.POINTAGE;
CREATE TABLE BTK_BI.GESTIONNAIRE   AS SELECT * FROM SYSTEM.GESTIONNAIRE;
\end{lstlisting}

\textbf{Pourquoi séparer ?} Un compte de restitution n'a pas à pouvoir écrire
dans les tables de production, et les requêtes d'un rapport ne doivent pas
peser sur l'application. C'est la même logique qui a mené à la chaîne ETL.

\subsection{La chaîne ETL}

\texttt{etl/etl\_agences.py} consolide les données opérationnelles en un
\textbf{datamart en étoile agrégé par agence}. Elle enchaîne cinq étapes :

\noindent\begin{tabularx}{\textwidth}{|>{\columncolor{btkpale}\bfseries}L{3.2cm}|X|}
\hline
Extract & lit \texttt{AGENCE}, \texttt{B\_UTILISATEURS}, \texttt{CLIENT\_BTK}, \texttt{B\_OBJECTIF}, \texttt{POINTAGE} --- par Oracle, par export CSV, ou depuis l'extrait agrégé livré avec le projet\\
\hline
Nettoyage & écarte les lignes sans \texttt{SK\_AGENCE}, type les colonnes\\
\hline
Transformation & agrège par agence : effectif, gestionnaires, clients, comptes, crédits, épargne, taux de présence\\
\hline
Contrôle & valeurs manquantes, valeurs négatives, agences en double --- le chargement est \textbf{interrompu} si un contrôle échoue\\
\hline
Load & écrit \texttt{dim\_agence}, \texttt{fait\_agence}, et le fichier consommé par la segmentation\\
\hline
\end{tabularx}

\vspace{0.6em}
Le c\oe{}ur de la transformation tient en quelques lignes de pandas :

\begin{lstlisting}[style=shell]
eff = employes.groupby("SK_AGENCE").agg(
        effectif=("SK_UTILISATEUR", "count"),
        nb_gestionnaires=("EST_GESTIONNAIRE", "sum")).reset_index()
cli = clients.groupby("SK_AGENCE").size().reset_index(name="nb_clients")
obj = objectifs.groupby("SK_AGENCE").agg(
        total_comptes=("COMPTES", "sum"),
        production_credits=("CREDITS", "sum"),
        collecte_epargne=("EPARGNE", "sum")).reset_index()

datamart = (agences.merge(eff, on="SK_AGENCE", how="left")
                   .merge(cli, on="SK_AGENCE", how="left")
                   .merge(obj, on="SK_AGENCE", how="left"))
\end{lstlisting}

Le résultat porte sur les \textbf{49 entités} du réseau extraites de la base.
Deux notebooks Jupyter rejouent la chaîne pas à pas
(\texttt{etl/etl\_btk.ipynb}) et la segmentation
(\texttt{clustering/segmentation\_btk.ipynb}) ; chacun se termine par une
cellule qui relance le script correspondant et \textbf{vérifie que les
résultats coïncident}.

\subsection{Power BI --- la construction du rapport}

\subsubsection{Connexion}

\texttt{Obtenir les données} $\rightarrow$ \textbf{Oracle} $\rightarrow$
serveur \texttt{localhost:1521/FREEPDB1}, compte \texttt{BTK\_BI}. Les vues
\texttt{V\_BI\_EMPLOYES}, \texttt{V\_BI\_CLIENTS}, \texttt{V\_BI\_OBJECTIFS} et
\texttt{V\_BI\_POINTAGE} sont importées.

\subsubsection{Thème et gabarit}

Le dossier \texttt{powerbi/} fournit \texttt{btk\_theme.json} (couleurs,
polices, styles des visuels, fond de visuel \texttt{\#212329}) et
\texttt{template\_background.png}, l'image de fond du canevas en
\textbf{1280$\times$720}. Le thème s'applique par
\texttt{Affichage $\rightarrow$ Thèmes}, le fond par
\texttt{Format de la page $\rightarrow$ Arrière-plan du canevas}.

\subsubsection{Les quatre pages}

\noindent\begin{tabularx}{\textwidth}{|>{\columncolor{btkpale}\bfseries}L{3.6cm}|X|}
\hline
1. Résumé & quatre cartes KPI, deux anneaux (rôles, présence), deux histogrammes (évolution des crédits, production par type), tableau des cinq premières agences\\
\hline
2. Clients & répartition du portefeuille : gouvernorat, secteur d'activité, type de client\\
\hline
3. Performance commerciale & comptes, crédits par nature, épargne, packs et cartes\\
\hline
4. Présence / RH & présents, retards, absences, taux de présence\\
\hline
\end{tabularx}

\vspace{0.5em}
Des segments communs --- \texttt{District}, \texttt{ANNEE\_MOIS},
\texttt{TYPE\_CLIENT} --- filtrent l'ensemble, avec un bouton de
réinitialisation. La barre latérale du fond sert de navigation entre les
quatre pages.

\subsubsection{Les mesures DAX}

Les indicateurs sont des \textbf{mesures}, calculées au moment de l'affichage
et donc sensibles aux filtres. Les plus représentatives :

\begin{lstlisting}[style=shell]
Nb Employes = DISTINCTCOUNT(B_UTILISATEURS[SK_UTILISATEUR])

Total Comptes =
      SUM(B_OBJECTIF[SOUSCRIPTION_COMPTE_CHEQUES_OA])
    + SUM(B_OBJECTIF[SOUSCRIPTION_COMPTE_EPARGNES_OA])
    + SUM(B_OBJECTIF[SOUSCRIPTION_COMPTE_COURANTS_OA])

Total Credits =
      SUM(B_OBJECTIF[PRODUCTION_CREDITS_CONSO_OA])
    + SUM(B_OBJECTIF[PRODUCTION_CREDITS_IMMO_OA])
    + SUM(B_OBJECTIF[PRODUCTION_CREDITS_INVESTISSEMENT_OA])

Taux Presence =
DIVIDE(
    CALCULATE(COUNTROWS(POINTAGE), POINTAGE[STATUT] IN {"PRESENT","RETARD"}),
    COUNTROWS(POINTAGE)
)
\end{lstlisting}

Le rôle, lui, est une \textbf{colonne calculée} et non une mesure : il
qualifie chaque ligne d'employé, il doit donc être matérialisé pour servir
d'axe d'analyse.

\begin{lstlisting}[style=shell]
Role =
SWITCH(
    TRUE(),
    B_UTILISATEURS[EST_GESTIONNAIRE] = 1, "Gestionnaire",
    B_UTILISATEURS[EST_GESTIONNAIRE] = 0 && B_UTILISATEURS[SK_AGENCE] = 31, "Hors Commercial",
    "Conseiller"
)
\end{lstlisting}

\begin{retenir}
\textbf{Mesure ou colonne calculée ?} C'est la question DAX la plus fréquente.
Une \textbf{colonne calculée} est évaluée une fois, à l'actualisation, et
occupe de la mémoire : elle sert quand on veut \emph{filtrer ou grouper} par
sa valeur. Une \textbf{mesure} est évaluée à l'affichage, dans le contexte des
filtres actifs : elle sert à \emph{agréger}. Ici, \texttt{Role} est une
colonne (on groupe par rôle), \texttt{Taux Présence} une mesure (il change
selon le mois ou l'agence sélectionnés).
\end{retenir}

\subsection{La segmentation des agences}

Le datamart alimente un apprentissage non supervisé
(\texttt{clustering/segmentation\_reelle.py}). Le siège est écarté --- 540
employés, il écraserait toute autre structure --- et les montants passent au
logarithme avant standardisation, leur distribution étant très asymétrique.

Quatre modèles sont comparés sur \textbf{la totalité} des agences, les points
qu'un modèle refuse de classer étant comptés comme un groupe à part entière,
pour qu'aucun ne soit avantagé par le fait d'avoir écarté les cas difficiles.

\noindent\begin{tabularx}{\textwidth}{|L{3.6cm}|c|c|X|}
\hline
\textbf{Modèle} & \textbf{Clusters} & \textbf{Silhouette} & \textbf{Lecture}\\
\hline
K-Means & 3 & 0,604 & \textbf{retenu}\\
\hline
DBSCAN & 4 & 0,517 & rejette CENTRALE comme bruit\\
\hline
GMM & 3 & 0,446 & \\
\hline
Agglomératif & 3 & 0,417 & \\
\hline
\end{tabularx}

\vspace{0.5em}
Trois segments interprétables en sortent : 34 agences d'exploitation,
6 entités de production sans portefeuille, 8 entités sans activité
enregistrée.
